\section{Introduction}
\label{sec:introduction}

This work was developed for the Investments \& Portfolio Management
competition, which asked us to manage a systematic long-short portfolio of
S\&P 500 stocks against the S\&P 500 Total Return benchmark. The central problem
is to choose a trading rule and assess whether its apparent outperformance
compensates for the risks and costs required to obtain it. Historical returns
alone cannot answer this question: data conventions, parameter selection,
leverage and execution all affect the result.

We address this problem through a shared research framework. Historical
prices, constituent membership, sectors and corporate actions feed
reproducible strategy searches and a portfolio ledger that records trades,
financing and execution costs. Applying the same accounting conventions to
momentum, reversal, low volatility and sector momentum allows us to compare
their net returns and risks.

Within this framework, we select momentum \#1292 from a five-candidate
historical shortlist and evaluate it separately over the competition window.
Risk measures and sector attribution explain the resulting performance,
reconciling excess return to allocation, pure stock selection, interaction and
account effects.
We also construct constrained efficient frontiers using our questionnaire-based
risk preference to study estimated return-risk trade-offs among the selected
stocks.

Section~\ref{sec:objective} sets out the investment objective and risk
calibration. Section~\ref{sec:data} describes the data and research framework,
and Section~\ref{sec:momentum} specifies the momentum rules.
Section~\ref{sec:research} presents historical selection and comparisons with
alternative strategies. Section~\ref{sec:competition} evaluates competition
performance and attribution, while Section~\ref{sec:frontier} examines efficient
frontiers. Finally, Section~\ref{sec:discussion} summarizes the findings and
priorities for further work.
